\section*{Independent review: capped occurrence propagation}
\textbf{Status: COMPUTATIONAL for the tests; exact kernel invariants below.}
The reviewer inspected the candidate separately and used a set-based simultaneous
unit-closure oracle, rather than the inherited sweep as the only oracle.
The reviewed domain consists of proper clauses, unique valid clause IDs in the
scope, and consistent masks over the declared variable range.

The independent run passed 1,532 partial-assignment closure and source-read
checks, 872 complete chain-order checks, three explicit empty-clause checks, and
734 observations of released index/queue containers at fallback entry.
Budgets include zero and one bit below, at, and above the calculated full-fit
threshold. Source literal reads are instrumented through clause iterators and
match acquisition plus event plus fallback-sweep counters. The run took less
than one second. These finite tests establish no general SAT complexity bound.

Every selected clause is initially queued. Thereafter every new assignment
queues all selected clauses containing the literal it falsifies, with duplicate
pending entries suppressed. Thus a newly unit or empty clause cannot become
unobserved. Each implication is sound, and exhaustion yields unit closure
unless a conflict is reported. Different conflict schedules need not return
identical partial assignments.

Let $s$ be the selected clause count, $w_i$ their original widths, and
$L=\sum_i w_i$. On successful acquisition, at most $s+L$ event probes occur:
one initial probe per clause and at most one trigger per indexed literal
occurrence. More specifically the literal-read bound, including acquisition,
is $\sum_i w_i(w_i+2)$. Counting zero-width clauses and queue/probe overhead gives
an abstract-operation bound $O(s+\sum_i w_i(w_i+2))$, subject to the stated
access and dictionary cost model. This is a phase bound, not the complete
solver bound. Fixed width gives $O(s+L)$ such operations regardless of clause
order. Unbounded-width clauses may be rescanned quadratically in their width.

The packed logical cap reserves the entire queue and membership vector before
building the occurrence table. Acquisition stops before the entry that would
exceed the cap is retained. All reads already performed, including the failed
entry's literal read, remain charged. Abort occurs only during construction,
whose active iterator refers to the original source clause, not an occurrence
list. Clearing the occurrence table, pending deque and membership set therefore
releases their entries before sweep fallback. The current frame's mask/scalar
scratch remains charged across that call. Ancestor descriptors, cache payload
and the inherited source locator are additional costs outside this cap.

The dictionaries, lists and deque are Python structures: the claimed cap is a
specified logical representation budget, not a Python heap or RSS bound.
Dictionary search, allocation and dense-integer bit operations are not given a
complete bit-machine time accounting. The source locator enables random access;
the sequential-replay counter does not simulate event probes on a tape.
A partial unsuccessful index adds overhead to fallback, and a full index may
use linear or greater encoded space. No favorable uniform time--space theorem
or closure of the P versus NP implication is asserted.
